\chapter{الكيوبتات، ومتجهات الحالة المركبة، والمرصودات}

\lectureoverview{يُعَدّ سبين إلكترون واحد أول كيوبت كمي بحق. ويمكن تحضيره بالنسبة إلى أي اتجاه مكاني، ومع ذلك لا يعطي كاشف ثابت إلا نتيجتين فقط. سنسمح للصورة المغناطيسية الكلاسيكية بأن تفشل، ثم نبني الفضاء المتجهي المركب، وقاعدة الاحتمال، ولغة المصفوفات اللازمة لوصف ما تفعله التجربة فعلاً.}

\section{سبين واحد ومعنيان لكلمة متجه}

للبت الكلاسيكي حالتان ممكنتان ولا شيء بينهما. ويمكن أن نطلق عليهما \(0\) و\(1\)، أو وجه العملة وظهرها، أو أعلى وأسفل؛ فهذه الأسماء مجرد علامات لنقطتين متميزتين في فضاء حالة كلاسيكي. والكيوبت الكمي يعطي أيضًا بديلين عند قياسه بطريقة ثابتة، لكنه لا يمكن وصفه بنقطتين فقط في فضاء الحالة. وأبسط مثال هو سبين الإلكترون.

يوجد في هذا المثال مفهومان مختلفان للمتجه، والخلط بينهما يجعل كل ما يلي غامضًا بلا داعٍ. الأول هو \emph{متجه مكاني} عادي: سهم في فضاء ثلاثي الأبعاد له مقدار واتجاه. وللإلكترون عزم مغناطيسي يتصرف، لأغراضنا الحالية، كمغناطيس صغير. ومقداره ثابت، في حين يمكن إسناد اتجاه مكاني إليه في إجراء تحضير.

والشيء الثاني هو متجه مجرد في فضاء متجهي مركب. وهذا \emph{متجه الحالة} سيصف الحالة الكمية. وليس هو ببساطة السهم المغناطيسي مكتوبًا بترميز آخر. إن العلاقة بين الاتجاهات المكانية ومتجهات الحالة علاقة حقيقية ومهمة، لكنها ليست تطابق هذين الشيئين. وسنصل إلى تلك العلاقة تدريجيًا.

\section{إخفاق الصورة الكلاسيكية}

لنبدأ عن عمد بقصة كلاسيكية. ضع مغناطيسًا صغيرًا في مجال مغناطيسي خارجي قوي. وما لم يكن مصطفًّا معه بالفعل، فإنه يقوم بحركة سبق حول المجال. ومع إشعاعه للطاقة، يتخامد هذا السبق ويستقر المغناطيس في اتجاهه ذي الطاقة الدنيا. وباختيار اتجاه المجال المطبَّق، وانتظار هذا الاسترخاء، ثم إزالة المجال، يبدو أننا نستطيع تحضير العزم المغناطيسي في أي اتجاه نريده.

ويقترح الجهاز نفسه طريقةً للكشف. شغّل مجالًا في اتجاه مختار وراقب الإشعاع المتولد بينما يستقر المغناطيس الصغير. ينبغي للمغناطيس الكلاسيكي المصطف مع المجال ألا يشع أي طاقة. أما إذا كان في البداية معاكسًا للمجال، فينبغي أن يشع المقدار الأعظمي، وينبغي لزاوية متوسطة أن تعطي مقدارًا متوسطًا. وفي تلك الصورة تتغير الطاقة المشعة تغيرًا مستمرًا مع الزاوية.

\begin{figure}[htbp]
\centering
\begin{coursefitfigure}
\begin{tikzpicture}[>=Latex,thick,scale=0.92]
  \draw[rounded corners] (-0.2,-0.2) rectangle (3.0,3.4);
  \node at (1.4,3.05) {\small preparation};
  \draw[->,very thick] (1.4,0.45) -- (1.4,2.55);
  \node[right] at (1.4,2.35) {\small applied field};
  \fill (1.4,1.35) circle (2.3pt);
  \draw[->,very thick] (1.4,1.35) -- (2.25,1.95);
  \node[below] at (1.4,0.15) {\small choose a direction and wait};

  \draw[->] (3.25,1.6) -- (4.25,1.6);

  \draw[rounded corners] (4.5,-0.2) rectangle (8.8,3.4);
  \node at (6.65,3.05) {\small measurement in a fixed direction};
  \draw[->,very thick] (5.15,0.45) -- (5.15,2.55);
  \fill (5.15,1.35) circle (2.3pt);
  \node[anchor=west] at (5.7,2.0) {no photon};
  \node[anchor=west] at (5.7,1.25) {one photon of fixed energy};
\node[anchor=west] at (5.7,0.5) {no intermediate outcome};
\end{tikzpicture}
\end{coursefitfigure}
\caption{اللغز العملي. يتبع تحضيرٌ متغير على نحو مستمر كاشفٌ لا يملك إلا خرجين ممكنين.}
\end{figure}

ذلك التنبؤ الكلاسيكي خاطئ في حالة الإلكترون. عند تشغيل مجال الكاشف لا تحدث إلا نتيجتان:
\begin{enumerate}
  \item لا ينبعث فوتون؛
  \item ينبعث فوتون واحد بطاقة ثابتة تفصل بين المستويين السفلي والعلوي.
\end{enumerate}
ولا يظهر قط فوتون يحمل كسرًا اعتباطيًا من تلك الطاقة. لقد انهارت الاستجابة الكلاسيكية المستمرة إلى جواب كمي ثنائي.

لم يصبح الاتجاه المستخدم في التحضير غير ذي صلة. فهو يضبط \emph{احتمالات} الإجابتين. فإذا حُضِّر الإلكترون إلى أعلى بالنسبة إلى الكاشف ظهرت نتيجة الأعلى بيقين. وإذا حُضِّر إلى أسفل وقع الانتقال بيقين. وإذا حُضِّر أفقيًا، فإن كاشفًا رأسيًا يعيد الإجابتين بالتواتر نفسه:
\begin{equation}
  P(\text{up})=\frac12,
  \qquad
  P(\text{down})=\frac12.
  \label{eq:l02-horizontal-probabilities}
\end{equation}
وهكذا تؤدي عائلة من التحضيرات المتغيرة باستمرار إلى قياس ذي قيمتين تتغير إحصاءاته تغيرًا مستمرًا.

\begin{classroomqa}
\lecturetimestamp{00:21:23}
\audiencequestion{هل يمكن لجسم صغير آخر، ربما إلكترون ثانٍ، أن يوفّر المجال المغناطيسي بدلًا من مغناطيس عياني؟}
\lecturerresponse{يمكن ذلك، لكن الإلكترونين عندئذٍ يشكلان نظامًا كميًا مشتركًا ويمكن أن يدخلا في حالة تشابك. وهذا بالضبط هو الاتجاه الذي يمضي إليه المقرر. نؤجل ذلك إلى أن نتمكن أولًا من فضاء حالة إلكترون واحد.}
\end{classroomqa}

\begin{classroomqa}
\lecturetimestamp{00:22:17}
\audiencequestion{هل لا بد من إيجاد توزيع الاحتمال بتكرار التحضير والقياس كاملين؟ ماذا يحدث إذا قسنا الإلكترون نفسه مرة أخرى؟}
\lecturerresponse{نعم، يتطلب التوزيع الأصلي تجارب متكررة مُحضَّرة على نحو متماثل. بعد القياس الأول، فإن تكرار القياس نفسه على ذلك الإلكترون يعطي النتيجة نفسها؛ إذ يكون القياس قد حضّره بالنسبة إلى ذلك الكاشف. وإن انتظار انتقال إشعاعي هو نفسه أحد إجراءات التحضير الممكنة.}
\end{classroomqa}

\begin{classroomqa}
\lecturetimestamp{00:23:48}
\audiencequestion{هل نتخيل إلكترونًا حرًا، وهل يُترك المجال المغناطيسي في وضع التشغيل باستمرار؟}
\lecturerresponse{نعزل درجة حرية السبين ونتجاهل موضع الإلكترون، سواء أكان يتحرك بحرية أم كان مقيَّدًا على نحو آخر. ويمكن تشغيل مغناطيس كهربائي كبير بسرعة من أجل التحضير أو الكشف. وقد صُممت التجربة المبسطة لكي تسأل عن السبين وحده.}
\end{classroomqa}

\begin{classroomqa}
\lecturetimestamp{00:25:22}
\audiencequestion{إذا حُضِّر الإلكترون تمامًا إلى أسفل وعُكس المجال، فهل ينبعث فوتون بيقين، ومن أين تأتي طاقته؟}
\lecturerresponse{في حالة المعاكسة التامة يحدث الانتقال باحتمال واحد. أما في تحضير مائل فيكون الاحتمال أصغر. وتأتي الطاقة في النهاية من الشغل المبذول في إنشاء المجال المغناطيسي الخارجي أو عكسه؛ فالجهاز جزء من حساب الطاقة.}
\end{classroomqa}

\begin{classroomqa}
\lecturetimestamp{00:27:31}
\audiencequestion{متى يظهر الفوتون؟ وهل الانتقال آني؟}
\lecturerresponse{زمن الانبعاث عشوائي، مثل زمن اضمحلال حالة غير مستقرة. ويمكن تحديد عمر مميز أو نصف عمر، ويقصّر المجال الأقوى ذلك المقياس. ويدخل عدم اليقين في الطاقة--الزمن في الوصف الأدق، لكن الوقائع العملية هنا هي المخرجان الممكنان وتكراراتهما المقاسة.}
\end{classroomqa}

\begin{classroomqa}
\lecturetimestamp{00:29:17}
\audiencequestion{هل ستنتج حزمة تحتوي مليون إلكترون مُحضَّرة تحضيرًا متماثلًا ببساطة شدةً ضوئية قابلة للقياس؟}
\lecturerresponse{نعم. تكشف شدة المجموعة الإحصائية عن متوسط احتمال الانتقال، بينما يبقى زمن انبعاث أي إلكترون بعينه عشوائيًا. وتُظهر المجموعة الإحصائية التوزيع الذي لا يمكن لمحاولة واحدة أن تُظهره.}
\end{classroomqa}

هذه هي نقطة الانطلاق. يمكن لنقطتين كلاسيكيتين منفصلتين أن ترمزا إجابتي الكاشف، لكنهما لا تستطيعان ترميز جميع التحضيرات التي تغيّر احتمالات تلك الإجابات. نحن بحاجة إلى فضاء حالة مختلف.

\section{الأعداد المركبة وهندسة فضاء الحالة}

فضاء الحالة المطلوب خطي: يمكن جمع عناصره وضربها بكميات عددية. وكمياته العددية هي أعداد مركبة. اكتب
\begin{equation}
  z=a+ib,
  \qquad
  i^2=-1,
  \qquad
  z^*=a-ib,
  \label{eq:l02-complex-number}
\end{equation}
حيث إن \(a\) و\(b\) حقيقيان. ويعطي أخذ المرافق المركب الكمية الموجبة
\begin{align}
  z^*z
  &=(a-ib)(a+ib) \notag\\
  &=a^2+b^2.
  \label{eq:l02-complex-norm}
\end{align}
وهذا الاختزال البسيط هو البذرة الجبرية لقاعدة الاحتمال.

\begin{figure}[htbp]
\centering
\includegraphics[width=0.74\linewidth]{lecture_02_figure_02.png}
\caption{تُستعرض على اللوح عملية أخذ المرافق المركب و\(z^*z\) قبل أن تُفسَّر السعات على أنها احتمالات.}
\lectureframe{00:37:10}
\end{figure}

يحتوي الفضاء المتجهي المركب على عناصر \(\ket a,\ket b,\ldots\) بحيث تكون
\begin{equation}
  c\ket a,
  \qquad
  \ket a+\ket b
\end{equation}
متجهاتٍ أيضًا متى كان \(c\) عددًا مركبًا. وتعطي متجهات الأعمدة تمثيلًا ملموسًا:
\begin{equation}
  \begin{aligned}
    \ket a&=
    \begin{pmatrix}
      a_1\\a_2\\ \vdots\\a_n
    \end{pmatrix},
    &
    c\ket a&=
    \begin{pmatrix}
      ca_1\\ca_2\\ \vdots\\ca_n
    \end{pmatrix},\\[0.7em]
    \ket a+\ket b&=
    \begin{pmatrix}
      a_1+b_1\\a_2+b_2\\ \vdots\\a_n+b_n
    \end{pmatrix}.
  \end{aligned}
  \label{eq:l02-vector-operations}
\end{equation}
بل إن الأعداد المركبة نفسها تُشكِّل فضاءً متجهيًا مركبًا أحادي البعد. أما لسبين الإلكترون فسنحتاج إلى بعدين.

نربط بكل كيت على هيئة عمود برا مرافقًا على هيئة صف:
\begin{equation}
  \ket a=
  \begin{pmatrix}a_1\\a_2\end{pmatrix}
  \quad\longleftrightarrow\quad
  \bra a=
  \begin{pmatrix}a_1^*&a_2^*\end{pmatrix}.
  \label{eq:l02-bra-ket-components}
\end{equation}
وعندئذٍ يكون الجداء الداخلي
\begin{equation}
  \braket{b}{a}=b_1^*a_1+b_2^*a_2,
  \label{eq:l02-inner-product}
\end{equation}
ويكون مربع المعيار
\begin{equation}
  \braket{a}{a}=|a_1|^2+|a_2|^2.
  \label{eq:l02-norm}
\end{equation}
وهو حقيقي وغير سالب لأن كل حد فيه من الصورة \(z^*z\).

\begin{classroomqa}
\lecturetimestamp{00:45:24}
\audiencequestion{لماذا لا تظهر نجمة صريحة على رمز البرا؟}
\lecturerresponse{إن تحويل الكيت إلى البرا الموافق له يعني أصلًا أخذ المنقول والمرافق المركب. فقلب اتجاه القوس هو الترميز لتلك العملية، لذا فإن كتابة نجمة أخرى ستكون تكرارًا لما هو مضمر.}
\end{classroomqa}

\section{أساس الكيوبت وقاعدة بورن}

لنعد الآن إلى الكاشف. ليدل \(\ket{+}\) على حالته العليا، و\(\ket{-}\) على حالته السفلى. وحالة سبين عامة هي تركيب خطي
\begin{equation}
  \ket\psi=a_+\ket{+}+a_-\ket{-}
  \quad\longleftrightarrow\quad
  \begin{pmatrix}a_+\\a_-\end{pmatrix}.
  \label{eq:l02-general-qubit}
\end{equation}
إشارتا الزائد والناقص داخل رمزي الكيت هما علامتان لا عمليتا حساب. وبمجرد اختيار هذا الأساس المرتب،
\begin{equation}
  \ket{+}=
  \begin{pmatrix}1\\0\end{pmatrix},
  \qquad
  \ket{-}=
  \begin{pmatrix}0\\1\end{pmatrix}.
  \label{eq:l02-basis-columns}
\end{equation}
وفي الواقع،
\begin{equation}
  a_+\begin{pmatrix}1\\0\end{pmatrix}
  +a_-\begin{pmatrix}0\\1\end{pmatrix}
  =
  \begin{pmatrix}a_+\\a_-\end{pmatrix}.
  \label{eq:l02-basis-expansion}
\end{equation}

\begin{figure}[htbp]
\centering
\includegraphics[width=0.72\linewidth]{lecture_02_figure_03.png}
\caption{تُؤخذ حالتا الكاشف بوصفهما عمودي الأساس القياسي.}
\lectureframe{00:58:10}
\end{figure}

المعاملات هي سعات احتمال. واحتمال كل ناتج من نواتج الكاشف هو مربع مقدار السعة المناظرة:
\begin{equation}
  P(+)=|a_+|^2,
  \qquad
  P(-)=|a_-|^2.
  \label{eq:l02-born-rule}
\end{equation}
ولأن هذه هي النواتج الوحيدة،
\begin{equation}
  \begin{gathered}
    |a_+|^2+|a_-|^2=1,\\
    \text{or equivalently}\qquad \braket{\psi}{\psi}=1.
  \end{gathered}
  \label{eq:l02-normalization}
\end{equation}
فإن متجه الحالة الفيزيائي يكون لذلك مطبَّعًا.

\begin{classroomqa}
\lecturetimestamp{00:56:54}
\audiencequestion{هل للعلامتَين زائد وناقص معنى جبري في المركبتين؟}
\lecturerresponse{هنا تُستعملان لوسم حالتي الأساس ومعاملاتهما. وكان يمكن أن نسمّيهما \(1\) و\(2\)، أو أعلى وأسفل. وعلى وجه الخصوص، فإن \(-\ket{+}\) مضاعف قياسي صحيح لـ \(\ket{+}\)؛ وليست هي الحالة \(\ket{-}\).}
\end{classroomqa}

إن ضرب حالة بعدد عقدي معياره واحد يترك جميع احتمالاتها بلا تغيير:
\begin{equation}
  \ket\psi\longmapsto e^{i\phi}\ket\psi,
  \qquad
  |e^{i\phi}a_\pm|^2=|a_\pm|^2.
  \label{eq:l02-global-phase}
\end{equation}
ولذلك فإن متجهين لا يختلفان إلا بطور كلي يمثّلان الحالة الفيزيائية نفسها. أمّا الطور \emph{النسبي} فمختلف. فاستبدال \(a_-\) وحده بـ \(-a_-\)، على سبيل المثال، لا يمكن عمومًا إزالته بضرب المتجه كله في طور واحد؛ إذ يمكنه أن يغيّر نتائج القياسات في أساس آخر.

\begin{classroomqa}
\lecturetimestamp{00:59:54}
\audiencequestion{هل \(\ket{+}\) و\(-\ket{+}\) و\(e^{i\phi}\ket{+}\) حالات مختلفة؟ وماذا عن تغيير إشارة أحد حدود التراكب؟}
\lecturerresponse{هي متجهات رياضية مختلفة، لكنها عندما يضرب العامل المتجه بأكمله ويكون معياره واحدًا تصف الحالة الفيزيائية نفسها. أمّا الإشارة أو الطور النسبي بين حدّي الأعلى والأسفل فيمكن رصده في قياسات مناسبة، ولذلك يغيّر الحالة.}
\end{classroomqa}

الأساس متعامد ومطبَّع:
\begin{align}
  \braket{+}{+}&=1,
  &\braket{-}{-}&=1,
  \notag\\
  \braket{+}{-}&=0,
  &\braket{-}{+}&=0.
  \label{eq:l02-orthonormality}
\end{align}
فعلى سبيل المثال،
\begin{equation}
  \braket{+}{-}
  =
  \begin{pmatrix}1&0\end{pmatrix}
  \begin{pmatrix}0\\1\end{pmatrix}
  =0.
\end{equation}

\begin{figure}[htbp]
\centering
\includegraphics[width=0.76\linewidth]{lecture_02_figure_04.png}
\caption{يُتحقَّق مباشرة على اللوح من الجداءات الداخلية الأربعة لأساس الأعلى/الأسفل.}
\lectureframe{01:05:00}
\end{figure}

تقدّم حالة متساوية الوزن أول تصحيح مفيد على اللوح. فالعمود
\begin{equation}
  \begin{pmatrix}1\\1\end{pmatrix}
\end{equation}
له معيار تربيعي \(2\) لا \(1\)، ولذلك فهو غير مُطبَّع. والقسمة على \(\sqrt2\) تصلحه:
\begin{equation}
  \ket{\rightarrow}
  =
  \frac{1}{\sqrt2}
  \begin{pmatrix}1\\1\end{pmatrix}
  =
  \frac{\ket{+}+\ket{-}}{\sqrt2}.
  \label{eq:l02-horizontal-state}
\end{equation}
وهو يعطي \(P(+)=P(-)=1/2\) في الكاشف الرأسي، ويقابل تحضيرًا أفقيًا واحدًا. وليس التصحيح من \((1,1)^T\) إلى \(2^{-1/2}(1,1)^T\) شكليًا: فالتطبيع هو ما يحوّل السعات المربعة إلى توزيع احتمالي كامل.

بعد الاستراحة يجدر بنا العودة إلى المعنيين لكلمة متجه. فالسهم المكاني له ثلاث مركبات حقيقية. أمّا متجه حالة كيوبت فله مركبتان عقديتان، أي أربعة أعداد حقيقية قبل فرض التطبيع واعتبار الطور الكلي غير مؤثِّر. والكائنان مرتبطان، لكن أحدهما ليس مجرد إعادة تسمية للآخر. ولسنا بحاجة بعدُ إلى الخريطة الهندسية الكاملة بينهما.

\begin{classroomqa}
\lecturetimestamp{01:10:56}
\audiencequestion{كيف نعرف بُعد فضاء الحالة لنظام كمي مختلف؟}
\lecturerresponse{تخبرنا التجربة بعدد النواتج المتنافية التي يستطيع قياس كامل أن يميّز بينها. فلفّ الإلكترون له ناتجان من هذا النوع، ومن ثم فإن فضاء حالته العقدي ثنائي البعد. أمّا نظام لفٍّ مغزلي مقداره واحد فله ثلاثة نواتج في القياس الموافق، ويتطلب ثلاث مركبات.}
\end{classroomqa}

\section{المرصودات الكلاسيكية ومتوسطات المجموعة الإحصائية}

بعد أن وصفنا الحالات، ننتقل إلى الكميات القابلة للقياس. في نظام كلاسيكي ذي حالات متقطعة \(n\)، يكون المرصود دالة حقيقية القيمة \(f_n\) على تلك الحالات. وإذا كانت الحالة غير يقينية وتوصف باحتمالات \(p_n\)، فإن متوسطها على المجموعة الإحصائية هو
\begin{equation}
  \langle f\rangle_{\mathrm{ens}}
  =\sum_n p_n f_n.
  \label{eq:l02-classical-expectation}
\end{equation}
ولا يعني هذا الترميز أن قياسًا منفردًا لا بد أن يساوي هذا العدد أصلًا.

بالنسبة إلى قطعة نقود، عيّن \(f_+=+1\) و\(f_-=-1\). وعندئذٍ
\begin{equation}
  \langle f\rangle_{\mathrm{ens}}
  =p_+-p_-.
  \label{eq:l02-coin-expectation}
\end{equation}
لقطعة نقود عادلة متوسط يساوي صفرًا. وإذا كان \(p_+=3/4\) و\(p_-=1/4\)، فالمتوسط هو \(1/2\). وليس أيٌّ من \(0\) أو \(1/2\) ناتجًا لرمية واحدة؛ بل إن كلًّا منهما يوجز كثيرًا من نتائج \(+1\) و\(-1\).

\begin{classroomqa}
\lecturetimestamp{01:23:30}
\audiencequestion{كيف يمكن أن تكون القيمة المتوقعة صفرًا أو نصفًا، مع أن أيا من القيمتين لا يمكن أن يظهر في تجربة واحدة؟}
\lecturerresponse{قد تكون عبارة القيمة المتوقعة مضلِّلة. فهي تعني متوسط مجموعة إحصائية كبيرة، لا القيمة التي ينبغي توقع رؤيتها في المحاولة التالية. فمتوسط متتالية عادلة من نتائج \(+1\) و\(-1\) يساوي صفرًا، واختلال بنسبة ثلاثة إلى واحد يعطي متوسطًا يساوي النصف.}
\end{classroomqa}

ومن المرصودات المفيدة على نحو خاص دالةُ المؤشر لحالة واحدة:
\begin{equation}
  \chi_k(n)=
  \begin{cases}
    1,&n=k,\\
    0,&n\ne k.
  \end{cases}
\end{equation}
ومتوسطها هو
\begin{equation}
  \langle\chi_k\rangle_{\mathrm{ens}}=p_k.
\end{equation}
وبالتالي فإن معرفة متوسطات جميع المرصودات الكلاسيكية تتضمن، على وجه الخصوص، معرفة التوزيع الاحتمالي.

\section{المصفوفات والمرصودات الكمومية}

في ميكانيكا الكم لا تُمثَّل المرصودات بدوال على نقاط فضاء الحالة، بل بمؤثرات خطية. وفي فضاء حالة منتهي البعد، تكون هذه المؤثرات مصفوفات. فالمصفوفة \(M\) تُحوِّل متجهًا إلى آخر:
\begin{equation}
  M\ket a=\ket{a'}.
  \label{eq:l02-matrix-action}
\end{equation}
وعند وجود مركبتين،
\begin{equation}
  \begin{pmatrix}
    m_{11}&m_{12}\\
    m_{21}&m_{22}
  \end{pmatrix}
  \begin{pmatrix}a_1\\a_2\end{pmatrix}
  =
  \begin{pmatrix}
    m_{11}a_1+m_{12}a_2\\
    m_{21}a_1+m_{22}a_2
  \end{pmatrix}.
  \label{eq:l02-matrix-components}
\end{equation}

\begin{figure}[htbp]
\centering
\includegraphics[width=0.78\linewidth]{lecture_02_figure_05.png}
\caption{يُفصَّل فعل المصفوفة صفًا صفًا إلى المركبتين اللتين يتألف منهما المتجه المتحوَّل.}
\lectureframe{01:32:10}
\end{figure}

تُظهر المتجهات الحقيقية ثنائية البعد الهندسة بوضوح. فالمصفوفات
\begin{equation}
  \begin{pmatrix}2&0\\0&2\end{pmatrix},
  \qquad
  \begin{pmatrix}1&0\\0&2\end{pmatrix},
  \qquad
  \begin{pmatrix}2&0\\0&1\end{pmatrix}
\end{equation}
تقوم، على الترتيب، بتمديد كل اتجاه على نحو منتظم، وتمديد المركبة \(y\) فقط، وتمديد المركبة \(x\) فقط. أمّا المصفوفة
\begin{equation}
  \begin{pmatrix}0&1\\1&0\end{pmatrix}
\end{equation}
فتبدّل بين المركبتين وتعكس المستوى حول \(x=y\). وفي الوقت نفسه،
\begin{equation}
  \begin{pmatrix}0&1\\-1&0\end{pmatrix}
  \begin{pmatrix}v_x\\v_y\end{pmatrix}
  =
  \begin{pmatrix}v_y\\-v_x\end{pmatrix}
  \label{eq:l02-quarter-turn}
\end{equation}
تدير المتجهات ربع دورة في اتجاه عقارب الساعة وفق الاصطلاح المستخدم على اللوح. ويعطي القص مثالًا خطيًا آخر، وإن كانت مصفوفتُه تُترك كتدريب. والقيد الجوهري هو الخطية: فليس كل تحويل يمكن تخيله للمستوى تُمثِّله مصفوفة.

تنتمي المرصودات الكمومية إلى فئة خاصة.

\begin{definition}
تكون المصفوفة \(M\) \emph{هيرميتية} إذا حققت
\begin{equation}
  M_{ij}=M_{ji}^*.
  \label{eq:l02-hermitian-condition}
\end{equation}
\end{definition}

يفرض الشرط القطري \(M_{ii}=M_{ii}^*\) أن يكون كل عنصر على القطر حقيقيًا. وتأتي العناصر خارج القطر في أزواج مترافقة. فعلى سبيل المثال،
\begin{equation}
  M=
  \begin{pmatrix}
    r&4+i\\
    4-i&s
  \end{pmatrix},
  \qquad r,s\in\mathbb R,
  \label{eq:l02-hermitian-example}
\end{equation}
مصفوفة هيرميتية. وإذا اتفق أن كانت جميع العناصر حقيقية، فإن كونها هيرميتية هو نفسه كونها متناظرة.

\begin{figure}[htbp]
\centering
\includegraphics[width=0.75\linewidth]{lecture_02_figure_06.png}
\caption{يُفرَد في نهاية المحاضرة شرطُ الترافق عبر القطر، \(M_{ij}=M_{ji}^*\).}
\lectureframe{01:42:50}
\end{figure}

ستمثّل المصفوفات الهيرميتية الكميات القابلة للقياس، بما في ذلك اللف المغزلي على امتداد الاتجاهات \(x\) و\(y\) و\(z\). ولا ينبغي الخلط بينها وبين المؤثرات التي تصف دفعة فيزيائية أو تطورًا زمنيًا.

\begin{classroomqa}
\lecturetimestamp{01:45:49}
\audiencequestion{إذا مثّلت مصفوفة دفعةً تغيّر الحالة إلى حالة أخرى، فهل يجب أن تحفظ التطبيع؟ وهل تفعل المصفوفة الهيرميتية ذلك؟}
\lecturerresponse{يجب على الدفعة الفيزيائية أن تحفظ طول متجه الحالة، لكن المصفوفة الهيرميتية العامة لا تفعل ذلك. والتحويلات التي تحفظ الطول وحدوية. تمثّل المؤثرات الهيرميتية المرصودات؛ وتمثّل المؤثرات الوحدوية التحويلات العكوسة التي تؤثر في الحالات. وتأتي علاقتهما التفصيلية بعد ذلك.}
\end{classroomqa}

لقد بلغنا الحد الأدنى من اللغة اللازمة للخطوة التالية. فالكيوبت متجه مُطبَّع في فضاء متجهي عقدي ثنائي البعد؛ والسعات تحدد الاحتمالات؛ والمرصودات مصفوفات هيرميتية. ومع اكتمال هذه العناصر، يمكن مناقشة اللف المغزلي على امتداد اتجاهات اعتباطية، ثم في نهاية المطاف الأزواج المتشابكة، من غير عودة إلى الصورة الكلاسيكية التي سبق أن فشلت.
